\documentclass{article}
\usepackage{amsmath,amssymb,amsthm}
\usepackage{leanify}
\newtheorem{lemma}{Lemma}
\begin{document}

A test of claims that follow a sentence which Leanify does not read. A claim is read
with the variables and hypotheses of the sentences of its result or proof that are read.
If a sentence before it is not read, the claim as read may not be the claim of the text.
The reports name every such claim, and a verified one is counted apart.

\section{A hypothesis that is not read}

The second sentence of the lemma is not read, so the claim lacks whatever it assumes. The
step of the proof is read without it as well.

\begin{lemma}\label{missing}
Let $x$ be a real number. Suppose moreover that the number is as large as the discussion requires. Then $x^2\ge0$.
\end{lemma}
\begin{proof}
We have $x^2\ge0$.
\end{proof}

\section{A case that has ended}

The third sentence of the proof turns to the other case and is not read. The last step is
therefore read under the hypothesis of the first case.

\begin{lemma}\label{cases}
For every real number $x$, $x\cdot x\ge0$.
\end{lemma}
\begin{proof}
Assume that $x\ge0$. Then $x\cdot x\ge0$. The other case is symmetric and we treat it now. Then $(-x)\cdot(-x)\ge0$.
\end{proof}

\section{A supposition that has ended}

The sentence that ends the argument by contradiction is not read. The conclusion of the
proof is read under the supposition, where it holds for an empty reason.

\begin{lemma}\label{contra}
For every real number $y$, $y^2+1>0$.
\end{lemma}
\begin{proof}
Suppose, for contradiction, that $y^2+1\le0$. Then $y^2\le-1$. But a square cannot be negative, which is absurd here. Hence $y^2+1>0$.
\end{proof}

\section{A proof that is read completely}

No claim of this lemma and its proof follows an unread sentence.

\begin{lemma}\label{complete}
For every real number $x$, $x\cdot x\ge0$.
\end{lemma}
\begin{proof}
Assume that $x\ge0$. Then $x\cdot x\ge0$. Assume now that $x<0$. Then $x\cdot x\ge0$.
\end{proof}

\end{document}
